% 分类目录：dl
\documentclass[12pt, a4paper]{article}
\usepackage[margin=2.5cm]{geometry}
\usepackage{ctex}
\usepackage{amsmath, amssymb}
\usepackage{listings}
\usepackage{xcolor}
\usepackage{tabularx}
\usepackage{hyperref}

\lstset{
    language=Python,
    basicstyle=\ttfamily\small,
    keywordstyle=\color{blue},
    commentstyle=\color{green!50!black},
    stringstyle=\color{red},
    showstringspaces=false,
    breaklines=true,
    numbers=left,
    numberstyle=\tiny\color{gray},
    frame=single
}

\title{知识点：参数初始化 (Parameter Initialization)}
\author{}
\date{}

\begin{document}
\maketitle

\section{核心定义}
\textbf{中文定义}：参数初始化是模型训练的第一步，指在训练开始前为神经网络的权重（Weights）和偏置（Biases）赋予初始值的过程。合适的初始化策略能够打破权重对称性，防止梯度消失或爆炸，从而加速模型的收敛。

\textbf{英文定义}：Parameter initialization is the process of assigning initial values to the weights and biases of a neural network before training. Proper initialization breaks symmetry, prevents gradient vanishing or exploding, and significantly accelerates convergence.

\section{为什么不能全 0 初始化？}
\textbf{打破对称性 (Symmetry Breaking)}：如果在一个隐藏层中所有神经元的权重均初始化为相同的值（如全 0 或相同常数），在前向传播时它们会得到相同的输出，在反向传播时也会得到相同的梯度更新。这意味着这层的所有神经元在训练过程中将始终保持一致，神经网络退化为一个线性模型，失去了多层特征提取的能力。

\section{经典初始化策略}

\subsection{随机初始化 (Random Initialization)}
通常采用高斯分布 $N(0, \sigma^2)$ 或均匀分布 $U(-k, k)$。
\\ \textbf{限制}：如果 $\sigma$ 过大，会导致激活值进入饱和区（梯度消失）；如果过小，则信号会逐层衰减。

\subsection{Xavier 初始化 (Glorot Initialization)}
\textbf{适用前提}：线性激活函数或 Tanh/Sigmoid（在其线性区）。
\\ \textbf{核心目标}：保持每一层输出的方差与输入的方差一致，且保持梯度的方差一致。
\\ \textbf{公式}：从均匀分布中采样：
$$ W \sim U\left( -\sqrt{\frac{6}{n_{in} + n_{out}}}, \sqrt{\frac{6}{n_{in} + n_{out}}} \right) $$
其中 $n_{in}$ 为输入神经元数量，$n_{out}$ 为输出神经元数量。

\subsection{He 初始化 (MSRA Initialization)}
\textbf{适用前提}：针对 ReLU 及其变体。
\\ \textbf{核心思想}：考虑到 ReLU 会将一半的输入置为 0，导致方差减半，因此需要将方差扩大 2 倍。
\\ \textbf{公式}：从高斯分布中采样：
$$ W \sim N\left( 0, \frac{2}{n_{in}} \right) $$

\section{数学推导：方差一致性原理}
考虑一层 $a = Wx + b$，假设组件独立且均值为 0，则：
$$ \text{Var}(a) = n \cdot \text{Var}(W) \text{Var}(x) $$
为了使输入的方差与输出一致（$\text{Var}(a) = \text{Var}(x)$），必须满足：
$$ \text{Var}(W) = \frac{1}{n_{in}} $$
这就是 Xavier 与 He 重大差异的理论来源：ReLU 激活函数的特殊非线性性质要求更高的权重方差。

\section{其他高级方案}
\begin{itemize}
    \item \textbf{偏置初始化}：通常初始化为 0。在使用 ReLU 时，有时初始化为 0.01 以减少神经元“坏死”。
    \item \textbf{正交初始化 (Orthogonal)}：有助于在非常深的网络或 RNN 中保持信号强度。
    \item \textbf{数据依赖型初始化}：如 Layer Normalization 或 BN 会在训练初期降低对初始化的敏感度。
\end{itemize}

\section{关键代码片段 (PyTorch)}
\begin{lstlisting}
import torch.nn as nn

class SimpleNet(nn.Module):
    def __init__(self):
        super(SimpleNet, self).__init__()
        self.fc1 = nn.Linear(784, 256)
        self.fc2 = nn.Linear(256, 10)
        
        # 1. Xavier 初始化 (适用 Tanh/Sigmoid)
        nn.init.xavier_uniform_(self.fc1.weight)
        
        # 2. He 初始化 (适用 ReLU)
        nn.init.kaiming_normal_(self.fc2.weight, mode='fan_in', nonlinearity='relu')
        
        # 3. 偏置常数初始化
        nn.init.constant_(self.fc1.bias, 0)

model = SimpleNet()
\end{lstlisting}

\section{深度面试题}

\textbf{问：Xavier 为什么不适合 ReLU 激活函数？}\\
答：Xavier 的推导假设激活函数是线性的且在原点对称。ReLU 将负信号直接截断，导致输出的均值不为 0 且方差减半。实验证明，如果对 ReLU 使用 Xavier，深层网络的输出会迅速萎缩至 0。

\textbf{问：初始化是否越随机越好？}\\
答：随机性是为了打破对称性，但方差必须受控。太随机（方差太大）会导致梯度爆炸，太拘束（方差太小）会导致梯度消失。理想的初始化是在信号逐层传递时保持其功率（能量）稳定。

\section{参考文献}
\begin{enumerate}
    \item Glorot, X., \& Bengio, Y. (2010). \textit{Understanding the Difficulty of Training Deep Feedforward Neural Networks}. AISTATS.
    \item He, K., et al. (2015). \textit{Delving Deep into Rectifiers: Surpassing Human-Level Performance on ImageNet Classification}. ICCV.
    \item Saxe, A. M., et al. (2013). \textit{Exact Solutions to the Nonlinear Dynamics of Learning in Deep Linear Networks}. ICLR.
\end{enumerate}

\end{document}
